\section{Discussion}
\label{sec:discussion}

\paragraph{Like-for-like caveat, quantified}
The expt\_gap result is not a like-for-like comparison of learning algorithms.
\begin{itemize}
\item \textbf{The leaders differ in the external information they use.} TPOT-Mat (steels) and MODNet (glass) used no external data: their inputs are composition descriptors computed from the formula, and they are fitted on the task's own training labels. Darwin, the leader on expt\_gap and expt\_is\_metal, fine-tunes the DARWIN natural-science large language model \citep{xie2023master,xie2023darwin}. Its pre-training corpus of scientific text is external knowledge. That text may include reported experimental band gaps or metallicity of some test compounds, and the extent of this overlap is unknown.
\item \textbf{Our pipeline uses external computed data.} It feeds PBE band gaps and formation energies from the Materials Project into the experts. It uses no external experimental labels.
\end{itemize}
The ablation (Section~\ref{sec:ablation}) puts a number on this.
\begin{itemize}
\item Removing the DFT block raises our band-gap MAE by \AblNoDftGapDelta~eV (\AblNoDftGapRel\%), which is \AblNoDftGapShareOfLeaderGain\% of our margin over Darwin.
\item The remaining DFT-free pipeline is statistically on par with Darwin (\AblNoDftGap{} versus \AtwoLeaderGap~eV, \AblNoDftGapVsLeaderPExpr), not ahead of it.
\item On expt\_is\_metal the DFT block likewise turns a result significantly below Darwin into parity.
\end{itemize}
Our DFT block and Darwin's pre-training are different kinds of external information, so neither comparison isolates the learning algorithm. The claim that survives is therefore:
\begin{itemize}
\item a DFT-free, fully nested stack is on par with the leader on expt\_gap, but below it on expt\_is\_metal;
\item disclosed DFT transfer adds a further, clearly significant improvement;
\item that improvement, together with the hurdle experts, takes the pipeline past the leader on expt\_gap and to parity on expt\_is\_metal.
\end{itemize}
The hurdle experts matter to a similar degree: without them, too, most of the margin disappears.

\paragraph{Relation to prior DFT-informed gap models}
That computed gaps are informative inputs or low-fidelity sources for experimental gaps is well established \citep{lee2016gap,pilania2017mf,kauwe2020bandgap,chen2021multifidelity,liu2022denoise,ko2024gap}.
\begin{itemize}
\item \textbf{\citet{lee2016gap}} showed, for 270 compounds and $G_0W_0$ targets, that combining Kohn--Sham gaps with elemental predictors in a nonlinear model reduces the RMSE from 0.59 to 0.24~eV, compared with a linear correction of the Kohn--Sham gap alone.
\item \textbf{\citet{chen2021multifidelity}} reported a 22--45\% lower experimental-gap MAE when low-fidelity PBE gaps were added to multi-fidelity graph networks.
\end{itemize}
Adding our DFT block lowers the MAE by \AblNoDftGapGainRel\% relative to the DFT-free variant. This is smaller than the improvement \citet{chen2021multifidelity} report, and it is obtained in a different setting:
\begin{itemize}
\item from a formula lookup, without any structure for the target compound;
\item on the official Matbench folds;
\item with the DFT block covering only 73\% of the compositions, the remainder being flagged as unmatched.
\end{itemize}
The results are not directly comparable, because the datasets, splits, targets and metrics differ.

\paragraph{Where the pipeline does not lead}
On glass, MODNet v0.1.12's public predictions reach a probability ROC-AUC of \AtwoLeaderAUCGlass, against our \AtwoAUCGlass{}. Neither the glass-formation descriptors nor any CPU-feasible expert we tried closed this gap.
On expt\_is\_metal and steels the pipeline is statistically on par with the leaders. On expt\_is\_metal this parity depends on the DFT block.
We did not audit the leaders' protocols, for example whether all of their tuning was confined to the training folds.

\paragraph{Unmerged leaderboard submissions}
The official leaderboard has not been updated since early 2024: the last merged pull request to the \texttt{matbench} repository dates from January 2024, and several benchmark submissions are open and unreviewed (checked 19 October 2026). Our ranks are therefore among merged entries only. Three open submissions bear directly on our results. Their numbers are as reported by their authors and have not yet been reviewed by the maintainers.
\begin{itemize}
\item \textbf{expt\_gap.} Zoey-MetalGate \citep{mbpr364} reports an MAE of 0.2775 $\pm$ 0.0051~eV with a composition-only, metal-classifier-gated gradient-boosting blend. This is prior art for our hurdle experts on this task. It uses no external data, so it is the natural reference for our DFT-free (\AblNoDftGap~eV) and hurdle-free (\AblNoHurdleGap~eV) ablation variants; we did not test these differences statistically. Our full pipeline (\AtwoFinalGap~eV) would remain first.
\item \textbf{steels.} Atomic-AffQRF \citep{mbpr371} reports about 75.8~MPa (mean over seven seeds). It combines extra-trees predictors on atomic fractions by out-of-fold, weighted-MAE linear-program stacking inside each training fold, a design closely related to ours. It would rank above our \AtwoFinalSteels~MPa and above the merged leader.
\item \textbf{glass.} A composition-only gradient-boosting submission \citep{mbpr363} reports a ROC-AUC of 0.8658, which does not change the ranking on this task.
\end{itemize}

\paragraph{Classification metrics}
The board's ``ROC-AUC'' for the classification leaders is, in effect, balanced accuracy. Darwin recorded labels. MODNet recorded probabilities that the \texttt{matbench} version used for its board entry thresholded at 0.5.
Our probability ROC-AUC on expt\_is\_metal (\AtwoAUCMetal) is therefore not comparable to the board's numbers.
We recommend reporting both quantities in separate columns.

\paragraph{The local correction}
The local residual correction that motivated this work transferred only weakly from chemistry to materials. Once its benefit had to survive a nested inner-CV check, it was never kept. We report it as a negative result.

\paragraph{Limitations}
\begin{itemize}
\item \textbf{Two looks at the test folds.} Run~2 was written after run~1 had been scored, so its motivation, though none of its choices, depended on the first test results. Both runs are reported.
\item \textbf{Candidate space from fold 0.} The candidate feature blocks and expert pool were developed on fold~0's training rows. Those rows contain the test rows of folds 1--4.
\item \textbf{Optimistic tuning OOF.} The tuned expert's OOF predictions are mildly optimistic, because tuning and stacking share inner splits. This can affect only the stack weights.
\item \textbf{Scope.} Four tasks, CPU-feasible experts only, five folds and a single seed limit the power and generality of the conclusions.
\item \textbf{Post-hoc ablation.} The ablation was designed after run-2 scoring, is descriptive and reports unadjusted $p$-values. Each component was removed one at a time, so the effects are not additive. It uses one seed, and does not separate the gap features from the formation-energy features within the DFT block.
\end{itemize}

\section{Conclusions}
\label{sec:conclusions}
Under two pre-registered, fully nested protocols on the Matbench v0.1 composition tasks, a cross-fitted stack of standard experts with a disclosed DFT feature block and a hurdle model for zero-inflated gaps:
\begin{itemize}
\item lowered the experimental band-gap MAE to \AtwoFinalGap~eV, below the best merged leaderboard entry (a fine-tuned scientific language model); without the DFT block the MAE was \AblNoDftGap~eV, on par with that leader;
\item was on par with the leaders for metallicity (only with the DFT block) and steel yield strength;
\item remained behind on glass formation.
\end{itemize}
The local residual correction that originally motivated the study did not survive nested validation.
